\begin{theorem}[Commutation of adjacent two-site basic support projectors]
    \label{thm:appendixB_two_site_basic_support_projectors_commute}
    \lean{MPSTensor.AppendixBStructuralData.twoSiteBasicSupportProjection_commute_lifts}
    \leanok
    \uses{thm:appendixB_virtual_bond_projectors_commute,
        thm:appendixB_physical_isometry_tensor_power,
        thm:appendixB_two_site_basic_support_projector}
    Let $(P_2)_{01}$ and $(P_2)_{12}$ be the two adjacent placements of the
    two-site basic support projector on three physical sites. Then
    $(P_2)_{01}(P_2)_{12} =(P_2)_{12}(P_2)_{01}$.
    This is the local three-site consequence of the parent-space construction
    and the basic-vector form in
    \cite[Definition~3.9, Equations~(3.16)--(3.18), and
        Appendix~B]{Cirac2016MPDO_arXiv}. It does not assert commutation for every
    chain length.
\end{theorem}

\begin{proof}
    \leanok
    The one-site map $U$ is an isometry. The preceding range identity
    therefore identifies the physical projector with the isometric transport
    of the virtual bond projector:
    $P_2 =U^{\otimes2}P_\varphi (U^*)^{\otimes2}$.
    On three sites, multiplication of the two adjacent transports cancels the
    middle factors by $U^*U=1$. Consequently their products are
    \begin{align}
         & U^{\otimes3}P_{01}P_{12}(U^*)^{\otimes3},
        \notag                                       \\
         & U^{\otimes3}P_{12}P_{01}(U^*)^{\otimes3},
        \notag
    \end{align}
    respectively. They agree by
    Theorem~\ref{thm:appendixB_virtual_bond_projectors_commute}.
    The following support-location schematic records the one-site embedding
    and the two overlapping physical supports. It does not depict the
    doubled $U\cdots U^*$ operator networks in the identities above.
    % Schematic formula: U_0--\varphi_{01}--U_1--\varphi_{12}--U_2, with
    % virtual supports around each \varphi bond and overlapping physical
    % supports around sites 01 and 12.
    % Ink-to-index: U_r maps its two virtual half-indices to the upper physical
    % index.  Each \varphi dot carries a rank-one virtual bond.  The inner
    % tinted regions are P_{01},P_{12}; the two overlapping outlines are the
    % transported physical supports (P_2)_{01},(P_2)_{12}.
    % Contractions: the four grid bonds are the internal virtual sums.
    % The two outer virtual and three upper physical indices remain free.
    % Boundary signature: (virtual-west, virtual-east, physical-up,
    % physical-down)=(1,1,3,0).  Regions are annotations, not contractions.
    % Source verdict: Equation (3.18) and II_RFP_MPS.png supply the local
    % isometric bond geometry.  Lines 1305--1307 state only that RFP implies a
    % nearest-neighbor commuting parent Hamiltonian immediately from the
    % structure; the explicit adjacent-projector transport is a blueprint
    % algebraic elaboration, and the schematic marks only its support geometry.
    \begin{center}
        \tenkzkernel
        \begin{tenkz}[rows={wire}, cols=5, west=open, east=open, physical=up]
            \tn[skin=box]{U_0} &
            \tn[physical=none, label pos=s]{\varphi_{01}} &
            \tn[skin=box]{U_1} &
            \tn[physical=none, label pos=s]{\varphi_{12}} &
            \tn[skin=box]{U_2}
            \tndeclare{species}{support}{hue=source:red}
            \tndeclare{species}{transported}{hue=source:blue}
            \tnmark[form=enclosure, species=support, tint, label pos=s]
                {(1,2)}{$\supp(P_{01})$}
            \tnmark[form=enclosure, species=support, tint, label pos=s]
                {(1,4)}{$\supp(P_{12})$}
            \tnmark[form=enclosure, species=transported, label pos=nw]
                {(1,1) .. (1,3)}{$\supp((P_2)_{01})$}
            \tnmark[form=enclosure, species=transported, label pos=ne]
                {(1,3) .. (1,5)}{$\supp((P_2)_{12})$}
        \end{tenkz}
    \end{center}
\end{proof}

\begin{definition}[Appendix B $AX$ coefficient representative]
    \label{def:appendixB_qax}
    \lean{MPSTensor.AppendixBStructuralData.appendixBQAXOnCoeffSpace}
    \leanok
    \uses{def:appendixB_two_site_basic_space, def:parent_interaction}
    The $AX$ coefficient-space representative is, before the $AX$-lift,
    the canonical two-site parent interaction of the core tensor,
    $\widehat Q_{AX}=q_2(A_{\mathrm{core}})$.
    The source projector $Q_{AX}$ of Definition~D.2 acts on
    $\mathcal H_A\otimes\mathcal H_X$. The displayed operator is the
    two-site coefficient representative that will be placed on the $AX$ face;
    it is not by itself the construction of that source projector.
\end{definition}

\begin{definition}[Appendix B $XB$ coefficient representative]
    \label{def:appendixB_qxb}
    \lean{MPSTensor.AppendixBStructuralData.appendixBQXBOnCoeffSpace}
    \leanok
    \uses{def:appendixB_qax}
    The $XB$ coefficient-space representative is, before the $XB$-lift, the
    same coefficient-space representative,
    $\widehat Q_{XB}=\widehat Q_{AX}$.
    The source projector $Q_{XB}$ of Definition~D.2 acts on
    $\mathcal H_X\otimes\mathcal H_B$. The displayed operator is the
    coefficient representative later placed on the $XB$ face.
\end{definition}

\begin{theorem}[The (AX) factor action]
    \label{thm:appendixB_qax_factor_action}
    \lean{MPSTensor.AppendixBStructuralData.appendixBQAX_factor_action}
    \leanok
    \uses{def:appendixB_qax, thm:pair_lift_kronecker_matrices}
    On the homogeneous three-single-site coefficient space
    $(\C^d)^{\otimes3}$, the hatted Appendix~B (AX) representative acts as
    $[\widehat Q_{AX}]_{AX} =[\widehat Q_{AX}]\otimes I_d$.
    This equation gives the coefficient-factor specialization of the left
    local action in \cite[Definition~D.2]{Cirac2016MPDO_arXiv}. It does not
    construct a projector for an arbitrary factorization
    $\mathcal H_A\otimes\mathcal H_X\otimes\mathcal H_B$.
\end{theorem}

\begin{proof}
    \leanok
    Apply the left identity in
    Theorem~\ref{thm:pair_lift_kronecker_matrices} to $\widehat Q_{AX}$.
\end{proof}

\begin{theorem}[The (XB) factor action]
    \label{thm:appendixB_qxb_factor_action}
    \lean{MPSTensor.AppendixBStructuralData.appendixBQXB_factor_action}
    \leanok
    \uses{def:appendixB_qxb, thm:pair_lift_kronecker_matrices}
    On the homogeneous three-single-site coefficient space
    $(\C^d)^{\otimes3}$, the hatted Appendix~B (XB) representative acts as
    $[\widehat Q_{XB}]_{XB} =I_d\otimes[\widehat Q_{XB}]$.
    This equation gives the coefficient-factor specialization of the right
    local action in \cite[Definition~D.2]{Cirac2016MPDO_arXiv}. It does not
    construct a projector for an arbitrary factorization
    $\mathcal H_A\otimes\mathcal H_X\otimes\mathcal H_B$.
\end{theorem}

\begin{proof}
    \leanok
    Apply the right identity in
    Theorem~\ref{thm:pair_lift_kronecker_matrices} to $\widehat Q_{XB}$.
\end{proof}

\begin{lemma}[Appendix B coefficient representatives and the parent interaction]
    \label{lem:appendixB_qax_qxb_parent_interaction}
    \lean{MPSTensor.AppendixBStructuralData.appendixBQAXOnCoeffSpace_eq_parentInteraction}
    \lean{MPSTensor.AppendixBStructuralData.appendixBQXBOnCoeffSpace_eq_parentInteraction}
    \leanok
    \uses{lem:appendixB_basis_change_parent_interaction, def:appendixB_qax,
        def:appendixB_qxb}
    As coefficient-space representatives,
    $\widehat Q_{AX}=q_2(A)$ and $\widehat Q_{XB}=q_2(A)$.
\end{lemma}

\begin{proof}
    \leanok
    Both operators are $q_2(A_{\mathrm{core}})$, and the Appendix~B basis
    change leaves $q_2$ unchanged.
\end{proof}

\begin{lemma}[Appendix B two-site representative idempotents]
    \label{lem:appendixB_qax_qxb_idempotent}
    \lean{MPSTensor.AppendixBStructuralData.appendixBQAXOnCoeffSpace_idempotent}
    \lean{MPSTensor.AppendixBStructuralData.appendixBQXBOnCoeffSpace_idempotent}
    \leanok
    \uses{lem:parent_interaction_idempotent,
        lem:appendixB_qax_qxb_parent_interaction}
    The two Appendix~B two-site coefficient representatives are idempotent:
    \begin{align}
        \widehat Q_{AX}^{2}
         & =\widehat Q_{AX},
        \notag               \\
        \widehat Q_{XB}^{2}
         & =\widehat Q_{XB}.
        \notag
    \end{align}
\end{lemma}

\begin{proof}
    \leanok
    By the preceding identification, both representatives are the canonical
    two-site parent interaction $q_2(A)$. The latter is an orthogonal
    projector, hence is idempotent.
\end{proof}

\begin{theorem}[Appendix B overlapping two-site commutation]
    \label{thm:appendixB_overlapping_two_site_commutation}
    \lean{MPSTensor.AppendixBStructuralData.hasOverlappingTwoSiteCommutation}
    \leanok
    \uses{thm:appendixB_two_site_basic_support_projectors_commute,
        thm:appendixB_two_site_basic_support_projector,
        lem:appendixB_qax_qxb_idempotent,
        def:overlapping_two_site_supports}
    The two Appendix~B coefficient representatives satisfy
    \begin{align}
        \widehat Q_{AX}^{\,2}
         & =\widehat Q_{AX},
        \notag               \\
        \widehat Q_{XB}^{\,2}
         & =\widehat Q_{XB},
        \notag
    \end{align}
    and their adjacent actions on the common three-site coefficient space
    commute:
    \begin{align}
        (\widehat Q_{AX})_{AX}(\widehat Q_{XB})_{XB}
         & =
        (\widehat Q_{XB})_{XB}(\widehat Q_{AX})_{AX}.
        \notag
    \end{align}
    Thus they satisfy the algebraic overlapping two-site condition of
    Definition~\ref{def:overlapping_two_site_supports}. This is the
    coefficient-space analog of the projector-and-commutator part of
    \cite[Definition~D.2]{Cirac2016MPDO_arXiv}. The hatted maps have not yet
    been identified with the source projectors on the explicitly factored
    spaces. By itself the result does not prove the kernel-intersection
    equation in that definition and makes no assertion for every chain length.
    The coefficient-space kernel equation is established in the next two
    results.
\end{theorem}

\begin{proof}
    \leanok
    Both coefficient representatives are the complement $1-P_2$ of the
    two-site basic support projector. Their idempotency was established in
    Lemma~\ref{lem:appendixB_qax_qxb_idempotent}. Since the two adjacent lifts
    of $P_2$ commute, their complementary lifts commute as well.
\end{proof}

\begin{theorem}[Three-site intersection of adjacent parent kernels]
    \label{thm:three_site_ground_space_adjacent_parent_kernels}
    \lean{MPSTensor.adjacent_localTerm_eq_zero_iff_mem_groundSpace_succ}
    \lean{MPSTensor.groundSpace_three_eq_adjacent_twoSite_parent_kernels}
    \leanok
    \uses{thm:local_term_two_three_ax_xb_lifts,
        thm:ground_space_intersection, def:parent_interaction, def:injective}
    If $A$ is injective, then
    $G_3(A) = \ker(q_2(A)_{AX}) \cap \ker(q_2(A)_{XB})$.
    More generally, on $L+1$ sites, simultaneous annihilation by the two
    adjacent length-$L$ parent interactions is equivalent to membership in
    $G_{L+1}(A)$ whenever $A$ is injective and $L>1$.
\end{theorem}

\begin{proof}
    \leanok
    By Theorem~\ref{thm:ground_space_intersection}, a three-site vector belongs
    to $G_3(A)$ exactly when its restrictions to both adjacent pairs belong to
    $G_2(A)$. Since $\ker q_2(A)=G_2(A)$, this is equivalent to annihilation
    by both adjacent terms. Theorem~\ref{thm:local_term_two_three_ax_xb_lifts}
    identifies these terms with the two displayed lifts.
\end{proof}

\begin{theorem}[Appendix B coefficient-space analog of the Appendix D kernel datum]
    \label{thm:appendixB_definition_d2_coefficient_analogue}
    \lean{MPSTensor.AppendixBStructuralData.hasAppendixD2ParentCommutingHamiltonian}
    \lean{MPSTensor.rfp_hasAppendixD2ParentCommutingHamiltonian_of_leftCanonical}
    \leanok
    \uses{thm:appendixB_overlapping_two_site_commutation,
        thm:three_site_ground_space_adjacent_parent_kernels,
        lem:appendixB_qax_qxb_parent_interaction}
    Let $A$ be injective and let an Appendix~B structural datum be fixed. Put
    \begin{align}
        \widehat K_{AXB}
         & =G_3(A),
        \notag      \\
        \widehat Q_{AX}
         & =q_2(A),
        \notag      \\
        \widehat Q_{XB}
         & =q_2(A),
        \notag
    \end{align}
    with the two copies of $q_2(A)$ lifted to the indicated faces of the
    three-site coefficient space. Then
    \begin{align}
        [(\widehat Q_{AX})_{AX},(\widehat Q_{XB})_{XB}]
         & =0,
        \notag \\
        \widehat K_{AXB}
         & =
        \ker((\widehat Q_{AX})_{AX})
        \cap\ker((\widehat Q_{XB})_{XB}).
        \notag
    \end{align}
    These are the coefficient-space analogs of the idempotence, commutation,
    and kernel-intersection equations in
    \cite[Definition~D.2]{Cirac2016MPDO_arXiv}. The hatted operators are not
    identified here with the source projectors on
    $\mathcal H_A\otimes\mathcal H_X$ and
    $\mathcal H_X\otimes\mathcal H_B$. For a normal left-canonical RFP tensor,
    injectivity follows from the RFP structural theorem and is not an
    additional hypothesis.
\end{theorem}

\begin{proof}
    \leanok
    Idempotency and commutation are
    Theorem~\ref{thm:appendixB_overlapping_two_site_commutation}. The two
    coefficient representatives are $q_2(A)$, and
    Theorem~\ref{thm:three_site_ground_space_adjacent_parent_kernels} supplies
    the remaining kernel equation.
    % Source-scope note: this Appendix D corollary is not the argument printed
    % for RFP => NNCPH.  The proof at source lines 1305--1307 instead refers
    % directly to the structural characterization.
\end{proof}

\begin{theorem}[Transport of the three-site commutator]
    \label{thm:three_site_commutator_chain_transport}
    \lean{MPSTensor.cyclicRestrictₗ_three_localTerm_left}
    \lean{MPSTensor.cyclicRestrictₗ_three_localTerm_right}
    \lean{MPSTensor.localTerm_adjacent_twoSite_commute_of_threeSite_zero_one_commute}
    \leanok
    \uses{def:local_term_parent}
    Suppose the first two length-two parent interactions commute on the
    three-site chain:
    $h_0^{(3)}h_1^{(3)} =h_1^{(3)}h_0^{(3)}$.
    Then, for every $N>2$ and every $i\in\mathbb Z/N\mathbb Z$, the two
    adjacent translates commute:
    $h_i^{(N)}h_{i+1}^{(N)} =h_{i+1}^{(N)}h_i^{(N)}$.
    For each exterior configuration $\tau$, let $R_{i,\tau}^{(3)}$ restrict to
    the cyclic three-site window beginning at $i$. It intertwines the two
    $N$-site terms with the two three-site terms:
    \begin{align}
        R_{i,\tau}^{(3)}h_i^{(N)}
         & =h_0^{(3)}R_{i,\tau}^{(3)},
        \notag                         \\
        R_{i,\tau}^{(3)}h_{i+1}^{(N)}
         & =h_1^{(3)}R_{i,\tau}^{(3)}.
        \notag
    \end{align}
    % Formula and contraction set: the two intertwining identities above are
    % the complete transport statement.  Restriction fixes every exterior
    % physical index to tau and leaves the three-site window free; there is no
    % tensor-network contraction beyond those algebraic evaluations.
    % Ink-to-index and boundary signature after demotion: there is no diagram
    % ink.  R_{i,tau}^{(3)} has an N-site input and a three-site output, while
    % h_i^{(N)},h_{i+1}^{(N)} and h_0^{(3)},h_1^{(3)} act on the spaces shown in
    % the displayed equations.
    % Source verdict: lines 1305--1307 only call RFP=>NNCPH immediate from the
    % structural characterization.  They do not state the explicit adjacent
    % projector or restriction lemmas used here; those are blueprint algebraic
    % elaborations.  The old ring schematic added unsupported source ink.
    Thus the same local $A\!X\!B$ calculation is placed at every cyclic
    translate through the two restriction identities above.
    The hypothesis $N>2$ is essential to this transport: at $N=2$ the two
    cyclic windows contain the same sites in opposite orders. This case
    is not part of \cite[Theorem~3.10]{Cirac2016MPDO_arXiv}.
\end{theorem}

\begin{proof}
    \leanok
    \uses{def:local_term_parent}
    Restrict an $N$-site vector to the cyclic window $(i,i+1,i+2)$ while
    holding all other physical indices fixed. Window extraction and
    replacement identify the two restricted local terms with
    $h_0^{(3)}$ and $h_1^{(3)}$. Apply the three-site commutator to the
    restricted vector and then evaluate at the original three-site
    configuration. Replacing a window by the values extracted from that
    configuration recovers the original $N$-site configuration.
\end{proof}

\begin{theorem}[Appendix B adjacent commutation on periodic chains]
    \label{thm:appendixB_adjacent_chain_commutation}
    \lean{MPSTensor.AppendixBStructuralData.adjacent_twoSite_localTerms_commute}
    \leanok
    \uses{thm:appendixB_overlapping_two_site_commutation,
        thm:three_site_commutator_chain_transport}
    For every Appendix~B structural datum and every periodic chain of length
    $N>2$, the adjacent translated length-two parent interactions commute.
\end{theorem}

\begin{proof}
    \leanok
    Theorem~\ref{thm:appendixB_overlapping_two_site_commutation} gives the
    three-site commutator, and
    Theorem~\ref{thm:three_site_commutator_chain_transport} places it at every
    adjacent pair on the periodic chain.
\end{proof}

\begin{definition}[Appendix B local-projector family]
    \label{def:appendixB_chain_local_projectors}
    \lean{MPSTensor.AppendixBStructuralData.hasProductPairLocalProjectors}
    \leanok
    \uses{thm:appendixB_adjacent_chain_commutation,
        lem:nncph_from_adjacent_two_site_commutation}
    For an Appendix~B structural datum and a periodic chain of length $N>2$,
    set $p_i=h_i^{(N)}$ for $i\in\mathbb Z/N\mathbb Z$, where $h_i^{(N)}$ is
    the translated length-two parent interaction. These maps are idempotent
    and commute pairwise: adjacent commutation follows from
    Theorem~\ref{thm:appendixB_adjacent_chain_commutation}, while disjoint
    length-two interactions commute by locality.
\end{definition}

\begin{theorem}[Left-canonical RFP commutation and zero energy]
    \label{thm:left_canonical_rfp_chain_commutation}
    \lean{MPSTensor.AppendixBStructuralData.isNNCPH}
    \lean{MPSTensor.rfp_implies_nncph_of_leftCanonical}
    \lean{MPSTensor.rfp_implies_nncph_ground_state_of_leftCanonical}
    \leanok
    \uses{def:appendixB_chain_local_projectors}
    A normal left-canonical renormalization fixed-point tensor of nonzero
    virtual dimension has pairwise
    commuting length-two parent interactions on every periodic chain of length
    $N>2$, and its periodic matrix product vector is annihilated by every such
    interaction. The removal of the left-canonical hypothesis and the
    ground-space spanning equation in
    \cite[Definition~3.9]{Cirac2016MPDO_arXiv} remain separate.
\end{theorem}

\begin{proof}
    \leanok
    The Appendix~B structural theorem supplies the datum used in
    Definition~\ref{def:appendixB_chain_local_projectors}. Its commuting
    idempotents are precisely the translated parent interactions. The
    standard frustration-free identity shows that each interaction annihilates
    the periodic matrix product vector.
\end{proof}

\begin{lemma}[Appendix B overlapping condition from the lifted commutator]
    \label{lem:appendixB_overlapping_condition_from_lift_commutator}
    \lean{MPSTensor.AppendixBStructuralData.hasOverlappingTwoSiteCommutation_of_commute_lifts}
    \leanok
    \uses{def:overlapping_two_site_supports, lem:appendixB_qax_qxb_idempotent}
    Suppose the lifted $AX$ and $XB$ representatives commute:
    \begin{align}
        (\widehat Q_{AX})_{AX}(\widehat Q_{XB})_{XB}
         & =
        (\widehat Q_{XB})_{XB}(\widehat Q_{AX})_{AX}.
        \notag
    \end{align}
    Then $\widehat Q_{AX}$ and $\widehat Q_{XB}$ satisfy the algebraic
    overlapping two-site condition of
    Definition~\ref{def:overlapping_two_site_supports}.
\end{lemma}

\begin{proof}
    \leanok
    The idempotency clauses are supplied by
    Lemma~\ref{lem:appendixB_qax_qxb_idempotent}. The only remaining clause
    in Definition~\ref{def:overlapping_two_site_supports} is exactly the
    displayed lifted commutator.
\end{proof}

\begin{lemma}[Definition D.2 projectors identified with Appendix B representatives]
    \label{lem:appendixB_d2_projectors_identified_with_representatives}
    \lean{MPSTensor.AppendixBStructuralData.hasOverlappingTwoSiteCommutation_of_appendixD2_identified_projectors}
    \lean{MPSTensor.AppendixBStructuralData.localTerm_two_three_zero_one_commute_of_appendixD2_identified_projectors}
    \leanok
    \uses{def:overlapping_two_site_supports, def:appendixB_qax,
        def:appendixB_qxb, lem:appendixB_overlapping_condition_parent_terms}
    Let $Q_{AX}$ and $Q_{XB}$ be the two-site coefficient-space transports of
    source projectors satisfying the Definition~D.2 local condition on
    $A X B$. If these transported projectors agree with the Appendix~B
    coefficient representatives,
    \begin{align}
        Q_{AX}
         & =\widehat Q_{AX},
        \notag               \\
        Q_{XB}
         & =\widehat Q_{XB},
        \notag
    \end{align}
    then $\widehat Q_{AX}$ and $\widehat Q_{XB}$ satisfy the algebraic
    overlapping two-site condition of
    Definition~\ref{def:overlapping_two_site_supports}. Consequently,
    $h_0^{(3)}(A,2)h_1^{(3)}(A,2) = h_1^{(3)}(A,2)h_0^{(3)}(A,2)$.
    % Maintainer note: this conditional identification statement does not
    % construct $Q_{AX}$ and $Q_{XB}$ from the Appendix~B basic-vector form.
\end{lemma}

\begin{proof}
    \leanok
    Rewriting the two projectors in the Definition~D.2 condition as
    $\widehat Q_{AX}$ and $\widehat Q_{XB}$ gives the algebraic overlapping
    condition. The three-site commutation is then
    Lemma~\ref{lem:appendixB_overlapping_condition_parent_terms}.
\end{proof}

\begin{lemma}[Three-site lifts of the Appendix B coefficient representatives]
    \label{lem:appendixB_three_site_parent_lifts}
    \lean{MPSTensor.AppendixBStructuralData.localTerm_two_three_zero_eq_appendixBQAXOnCoeffSpace_lift}
    \lean{MPSTensor.AppendixBStructuralData.localTerm_two_three_one_eq_appendixBQXBOnCoeffSpace_lift}
    \leanok
    \uses{thm:local_term_two_three_ax_xb_lifts,
        lem:appendixB_qax_qxb_parent_interaction}
    On the three-site $A,X,B$ window, the two adjacent length-two local terms
    are the two local lifts
    \begin{align}
        h_0^{(3)}(A,2)
         & =(\widehat Q_{AX})_{AX},
        \notag                      \\
        h_1^{(3)}(A,2)
         & =(\widehat Q_{XB})_{XB}.
        \notag
    \end{align}
\end{lemma}

\begin{proof}
    \leanok
    Apply the three-site local-term identities to $q_2(A)$, and
    replace $q_2(A)$ by $\widehat Q_{AX}$ or $\widehat Q_{XB}$ using the
    preceding two-site identifications.
\end{proof}

\begin{lemma}[Appendix B overlapping condition gives three-site local-term commutation]
    \label{lem:appendixB_overlapping_condition_parent_terms}
    \lean{MPSTensor.AppendixBStructuralData.localTerm_two_three_zero_one_commute_of_overlapping}
    \leanok
    \uses{lem:appendixB_three_site_parent_lifts}
    If the Appendix~B $AX$ and $XB$ coefficient representatives satisfy the
    overlapping two-site commutation datum
    \begin{align}
        \widehat Q_{AX}^2
         & =\widehat Q_{AX},
        \notag               \\
        \widehat Q_{XB}^2
         & =\widehat Q_{XB},
        \notag               \\
        (\widehat Q_{AX})_{AX}(\widehat Q_{XB})_{XB}
         & =
        (\widehat Q_{XB})_{XB}(\widehat Q_{AX})_{AX},
        \notag
    \end{align}
    then the corresponding first two translated length-two local terms on the
    three-site window commute:
    $h_0^{(3)}(A,2)h_1^{(3)}(A,2) = h_1^{(3)}(A,2)h_0^{(3)}(A,2)$.
\end{lemma}

\begin{proof}
    \leanok
    By Lemma~\ref{lem:appendixB_three_site_parent_lifts}, the two translated
    local terms are exactly the $AX$ and $XB$ lifts of
    $\widehat Q_{AX}$ and $\widehat Q_{XB}$. The displayed equality is
    therefore the commutation clause of the overlapping two-site datum.
\end{proof}

\begin{lemma}[Appendix B lifted commutator gives three-site local-term commutation]
    \label{lem:appendixB_lift_commutator_parent_terms}
    \lean{MPSTensor.AppendixBStructuralData.localTerm_two_three_zero_one_commute_of_commute_lifts}
    \lean{MPSTensor.AppendixBStructuralData.localTerm_two_three_zero_one_commute_of_appendixD2}
    \leanok
    \uses{lem:appendixB_overlapping_condition_from_lift_commutator,
        lem:appendixB_overlapping_condition_parent_terms}
    If the Appendix~B $AX$ and $XB$ coefficient representatives satisfy the
    lifted Definition~D.2 commutator
    \begin{align}
        (\widehat Q_{AX})_{AX}(\widehat Q_{XB})_{XB}
         & =
        (\widehat Q_{XB})_{XB}(\widehat Q_{AX})_{AX},
        \notag
    \end{align}
    then the corresponding first two translated length-two local terms on the
    three-site window commute:
    $h_0^{(3)}(A,2)h_1^{(3)}(A,2) = h_1^{(3)}(A,2)h_0^{(3)}(A,2)$.
    The same conclusion holds if one supplies a Definition~D.2-shaped local
    condition whose two operators are these coefficient-space representatives,
    rather than the source projectors themselves.
\end{lemma}

\begin{proof}
    \leanok
    The lifted commutator, together with the Appendix~B idempotency already
    recorded above, gives the overlapping two-site commutation datum. A supplied
    Definition~D.2-shaped local condition on the coefficient representatives
    implies the same conclusion by retaining its commutation clause and not using
    the kernel-intersection condition
    \begin{align}
        K_{AXB}
         & =
        \ker((\widehat Q_{AX})_{AX})
        \cap
        \ker((\widehat Q_{XB})_{XB}).
        \notag
    \end{align}
\end{proof}
